% BHU PYQ -- Manually transcribed & error-corrected
% Paper: BCH-224 : Business Mathematics
% Exam: B.Com. (Hons.) (Semester IV) Examination, 2015-16

\documentclass[11pt,a4paper]{article}
\usepackage[a4paper,top=2.5cm,bottom=2.5cm,left=2.8cm,right=2.8cm,headheight=14pt]{geometry}
\usepackage{amsmath,amssymb}
\usepackage[T1]{fontenc}
\usepackage[utf8]{inputenc}
\usepackage{lmodern,microtype,fancyhdr,enumitem,hyperref,lastpage,parskip}

\hypersetup{
    colorlinks=true,
    urlcolor=black,
    linkcolor=black,
    pdftitle={BCH-224: Business Mathematics -- BHU B.Com. (Hons.) Sem IV 2015-16}
}

\pagestyle{fancy}
\fancyhf{}
\renewcommand{\headrulewidth}{0.4pt}
\renewcommand{\footrulewidth}{0.4pt}
\fancyhead[L]{\small   \textit{Banaras Hindu University}}
\fancyhead[R]{\small   \textit{BCH-224: Business Mathematics}}
\fancyfoot[C]{\small Page  \thepage\ of \pageref{LastPage}}


\newlist{parts}{enumerate}{2}
\setlist[parts,1]{label=(\alph*),leftmargin=*,itemsep=5pt,topsep=3pt}
\setlist[parts,2]{label=(\roman*),leftmargin=*,itemsep=3pt,topsep=2pt}

\newcommand{\pts}[1]{\hfill   \textit{[#1]}}


\begin{document}


\begin{center}
  {\large\bfseries BANARAS HINDU UNIVERSITY}\\[2pt]
  {
\normalsize B.Com. (Hons.) --- Semester IV Examination, 2015-16}\\[6pt]
  \rule{0.8\linewidth}{0.5pt}\\[6pt]
  {\large\bfseries Paper No. BCH-224: Business Mathematics}\\[4pt]
  {
\normalsize Time: 3 Hours\hfill Maximum Marks: 70}
\end{center}

\rule{\linewidth}{0.4pt}

\smallskip



\noindent   \textit{Instructions: Answer all questions. All questions carry equal marks (14 marks each).}

\bigskip




\noindent   \textbf{Question 1.}

\begin{parts}
  \item Find the sum of the first $n$ odd natural numbers.\pts{7}
  \item If $a, b, c$ are in A.P., then prove that:
  \[
  \frac{1}{\sqrt{b}+\sqrt{c}}, \, \frac{1}{\sqrt{c}+\sqrt{a}}, \, \frac{1}{\sqrt{a}+\sqrt{b}}
  \]
  are also in A.P.\pts{7}
\end{parts}

\centerline{   \textbf{OR}}


\begin{parts}
  \item The $p^{ \text{th}}$ and $q^{ \text{th}}$ terms of an A.P. are $c$ and $d$ respectively. Find the $r^{ \text{th}}$ term.\pts{7}
  \item The sum of $n$ terms of an A.P. is 40, the common difference is 2, and the last term is 13. Find the number of terms.\pts{7}
\end{parts}

\medskip\rule{\linewidth}{0.2pt}




\noindent   \textbf{Question 2.}

\begin{parts}
  \item The sum of three numbers in A.P. is 21. If the second number is reduced by 1 and the third number is increased by 1, then the numbers obtained are in G.P. Find the numbers.\pts{7}
  \item If $a, b, c$ are in G.P. and $x, y$ are the arithmetic means between $a, b$ and $b, c$ respectively, prove that:
  \[
  \frac{a}{x} + \frac{c}{y} = 2
  \]\pts{7}
\end{parts}

\centerline{   \textbf{OR}}


\begin{parts}
  \item If the arithmetic mean of two numbers is 10 and their geometric mean is 6, then find the numbers.\pts{7}
  \item If $\frac{1}{x+y}, \frac{1}{2y}, \frac{1}{y+z}$ are in A.P., then prove that $x, y, z$ are in G.P.\pts{7}
\end{parts}

\medskip\rule{\linewidth}{0.2pt}

\pagebreak




\noindent   \textbf{Question 3.}

\begin{parts}
  \item If ${}^{2n+1}P_{n-1} : {}^{2n-1}P_n = 3:5$, then find the value of $n$.\pts{7}
  \item How many numbers between 100 and 10,000 can be made with the 8 digits 1, 2, 3, 0, 4, 5, 6, 7? (Repetition of digits is not allowed.)\pts{7}
\end{parts}

\centerline{   \textbf{OR}}


\begin{parts}
  \item How many triangles can be formed by joining the angular points of a decagon? How many diagonals does it have?\pts{7}
  \item A box contains 7 red, 6 white, and 4 blue balls. Find how many selections of three balls can be made so that none of them is red.\pts{7}
\end{parts}

\medskip\rule{\linewidth}{0.2pt}




\noindent   \textbf{Question 4.}

\begin{parts}
  \item Find the term independent of $x$ in the expansion of $\left(x^2 + \frac{1}{x}\right)^{12}$.\pts{7}
  \item In the expansion of $(1 + x)^{43}$, the coefficients of the $(2r + 1)^{ \text{th}}$ term and $(r + 2)^{ \text{th}}$ term are equal. Find the value of $r$.\pts{7}
\end{parts}

\centerline{   \textbf{OR}}


\begin{parts}
  \item Find the $19^{ \text{th}}$ term of $\left(2x^{1/2} - y^{1/3}\right)^{20}$.\pts{7}
  \item Find the coefficient of $x^{15}$ in the expansion of $(x - x^2)^{10}$.\pts{7}
\end{parts}

\medskip\rule{\linewidth}{0.2pt}

\pagebreak




\noindent   \textbf{Question 5.}

\begin{parts}
  \item If $f(x) = \log\left(\frac{1-x}{1+x}\right)$, then prove that:
  \[
  f(a) + f(b) = f\left(\frac{a+b}{1+ab}\right)
  \]\pts{4}
  \item Evaluate:
  \[
  \lim_{x  o 0} \frac{\sqrt{1+x} - \sqrt{1-x}}{x}
  \]\pts{4}
  \item If $y = x^{x^{x^{\dots \infty}}}$, then prove that:
  \[
  x \frac{dy}{dx} = \frac{y^2}{1 - y\log x}
  \]\pts{6}
\end{parts}

\centerline{   \textbf{OR}}


\begin{parts}
  \item Simplify:
  \[
  \int \left(\sqrt{x} + \frac{1}{\sqrt{x}}\right)^2 \, dx
  \]\pts{4}
  \item Simplify:
  \[
  \int \frac{x}{\sqrt{x+1}} \, dx
  \]\pts{4}
  \item Simplify:
  \[
  \int \frac{x+1}{(x+2)^2} \, dx
  \]\pts{6}
\end{parts}

\vfill
\rule{\linewidth}{0.4pt}

\begin{center}\small
\href{https://bhu-exam.vercel.app/}{Click here to visit our website}\\[4pt]
\href{https://me-aryan.vercel.app/}{Click here to visit our contributor}\\[4pt]
\href{https://quashan-me.vercel.app/}{Click here to visit our contributor}
\end{center}

\end{document}
